\section{Functions}

This section establishes standards for function design, implementation, and documentation in the STAR firmware codebase.

\subsection{Return Value Conventions}

\begin{itemize}
    \item If the name of a function is an action (or an imperative command), it should return an \textbf{error-code integer} (\texttt{esp\_err\_t}).
    \item If the name is a predicate (a function that returns true/false), it should return a \textbf{succeeded boolean}.
\end{itemize}

\begin{lstlisting}[language=C]
/* Action functions return esp_err_t */
esp_err_t star_bus_manager_init(...);
esp_err_t star_bus_manager_add_bus(...);
esp_err_t star_bus_config_destroy(...);

/* Predicate functions return bool */
bool error_handler_can_retry(error_handler_t* handler);
bool is_bus_initialized(const star_bus_config_t* config);

/* Getter functions return the value or pointer */
star_bus_config_t* star_bus_manager_find_bus(...);
const char* star_bus_type_to_string(star_bus_type_t type);
\end{lstlisting}

\subsection{Private and Exposed Private Functions}

\subsubsection{Private Functions (static)}

Functions internal to a single source file should be declared \texttt{static} with the \texttt{internal\_} prefix.

\begin{lstlisting}[language=C]
/* In .c file - private static functions */
static esp_err_t internal_register_bus_pin(star_pin_interface_t* pin_iface,
                                            gpio_num_t            pin,
                                            const char*           bus_name,
                                            const char*           usage,
                                            bool                  can_be_shared)
{
    if (pin < 0 || pin >= GPIO_NUM_MAX) {
        return ESP_OK;  /* Not an error, pin is not used */
    }
    /* ... implementation ... */
}

static esp_err_t internal_i2c_execute_with_retry(const star_bus_config_t* config,
                                                  i2c_port_t               port,
                                                  i2c_cmd_handle_t         cmd,
                                                  const char*              operation_name)
{
    /* ... implementation ... */
}

/* Declare prototypes at top of file */
static esp_err_t internal_register_bus_pin(...);
static esp_err_t internal_unregister_bus_pin(...);
static esp_err_t internal_register_config_pins(...);
\end{lstlisting}

\subsubsection{Exposed Private Functions (priv\_)}

If a function cannot be \texttt{static} and must be exposed across files within a module (but not publicly), prefix its name with \texttt{priv\_}.

\begin{lstlisting}[language=C]
/* In internal header (e.g., star_bus_internal.h) */
esp_err_t priv_star_bus_validate_config(const star_bus_config_t* config);
esp_err_t priv_star_bus_acquire_mutex(star_bus_manager_t* manager);
void      priv_star_bus_release_mutex(star_bus_manager_t* manager);
\end{lstlisting}

\subsection{Parameter Validation}

Validate all parameters at the start of public functions.

\begin{lstlisting}[language=C]
esp_err_t star_bus_manager_add_bus(star_bus_manager_t* manager,
                                    star_bus_config_t*  config)
{
    /* Validate required parameters */
    ESP_RETURN_ON_FALSE(manager && config,
                        ESP_ERR_INVALID_ARG,
                        s_TAG,
                        "Manager or config pointer is NULL");

    /* Validate config contents */
    ESP_RETURN_ON_FALSE(config->name && strlen(config->name) > 0,
                        ESP_ERR_INVALID_ARG,
                        manager->tag,
                        "Config name is NULL or empty");

    /* Validate enum range */
    ESP_RETURN_ON_FALSE(config->type > k_star_bus_type_none &&
                        config->type < k_star_bus_type_count,
                        ESP_ERR_INVALID_ARG,
                        manager->tag,
                        "Invalid bus type in config");

    /* Validate state */
    ESP_RETURN_ON_FALSE(manager->mutex,
                        ESP_ERR_INVALID_STATE,
                        manager->tag,
                        "Manager mutex not initialized");

    /* ... rest of implementation ... */
}

esp_err_t star_bus_i2c_write(const star_bus_manager_t* manager,
                              const char*               name,
                              const uint8_t*            data,
                              size_t                    len,
                              uint8_t                   reg_addr,
                              size_t*                   bytes_written)
{
    /* Validate all required pointers */
    ESP_RETURN_ON_FALSE(manager && name && data,
                        ESP_ERR_INVALID_ARG,
                        s_TAG,
                        "Manager, name, or data is NULL");

    /* Validate numeric constraints */
    ESP_RETURN_ON_FALSE(len > 0, ESP_ERR_INVALID_ARG, s_TAG,
                        "Write length must be > 0");

    /* Clear output parameters initially */
    if (bytes_written) {
        *bytes_written = 0;
    }

    /* ... rest of implementation ... */
}
\end{lstlisting}

\subsection{Function Declaration Formatting}

Align parameters vertically when they span multiple lines.

\begin{lstlisting}[language=C]
/* CORRECT - Aligned parameters */
esp_err_t star_bus_manager_init(star_bus_manager_t*     manager,
                                const char*             tag,
                                star_error_interface_t* error_iface,
                                star_pin_interface_t*   pin_iface);

star_bus_config_t* star_bus_config_create_spi_device(
    const char*                          name,
    spi_host_device_t                    host,
    gpio_num_t                           copi_pin,
    gpio_num_t                           cipo_pin,
    gpio_num_t                           sclk_pin,
    int32_t                              dma_chan,
    const spi_device_interface_config_t* dev_cfg);

/* Function pointers in structs */
typedef struct star_error_interface {
    esp_err_t (*record_error)(void*       ctx,
                              esp_err_t   error,
                              const char* message,
                              const char* file,
                              int         line,
                              const char* func);
    bool (*can_retry)(void* ctx);
    esp_err_t (*reset_state)(void* ctx);
    void* ctx;
} star_error_interface_t;
\end{lstlisting}

\subsection{Doxygen Documentation}

All public functions must have complete Doxygen documentation in header files.

\begin{lstlisting}[language=C]
/**
 * @brief Add a new bus/device configuration to the manager.
 *
 * Adds a previously created configuration (via star_bus_config_create_*)
 * to the manager's internal list. The manager takes ownership conceptually,
 * but the configuration still needs to be destroyed eventually (typically via
 * star_bus_manager_remove_bus or star_bus_manager_deinit).
 *
 * The configuration must have a unique name. This function is thread-safe.
 *
 * @param[in] manager Pointer to the initialized bus manager.
 *                    Must not be NULL.
 * @param[in] config  Pointer to the bus configuration structure to add.
 *                    Must not be NULL and must have a valid name.
 * @return esp_err_t ESP_OK on success,
 *                   ESP_ERR_INVALID_ARG if manager or config is invalid,
 *                   ESP_ERR_INVALID_STATE if a bus with the same name
 *                       already exists or mutex error,
 *                   ESP_ERR_TIMEOUT if mutex times out.
 */
esp_err_t star_bus_manager_add_bus(star_bus_manager_t* manager,
                                    star_bus_config_t*  config);
\end{lstlisting}

\subsection{Private Function Documentation}

Static functions should have brief documentation in the source file.

\begin{lstlisting}[language=C]
/**
 * @brief Helper: Execute I2C command with retry logic
 *
 * Wraps i2c_master_cmd_begin() with automatic retry for transient errors.
 *
 * @param[in] config         Pointer to I2C bus configuration
 * @param[in] port           I2C port number
 * @param[in] cmd            I2C command handle
 * @param[in] operation_name Name of operation for logging
 * @return esp_err_t Final result after retries
 */
static esp_err_t internal_i2c_execute_with_retry(const star_bus_config_t* config,
                                                  i2c_port_t               port,
                                                  i2c_cmd_handle_t         cmd,
                                                  const char*              operation_name);
\end{lstlisting}

\subsection{Function Implementation Patterns}

\subsubsection{Single Exit Point with Cleanup}

\begin{lstlisting}[language=C]
esp_err_t star_bus_manager_remove_bus(star_bus_manager_t* manager,
                                       const char*         name)
{
    ESP_RETURN_ON_FALSE(manager && name && strlen(name) > 0,
                        ESP_ERR_INVALID_ARG, s_TAG, "Invalid arguments");

    star_bus_config_t* to_remove = NULL;
    star_bus_config_t* prev      = NULL;
    esp_err_t          ret       = ESP_ERR_NOT_FOUND;

    /* Acquire mutex */
    if (xSemaphoreTake(manager->mutex,
            pdMS_TO_TICKS(CONFIG_STAR_KCONFIG_BUS_MUTEX_TIMEOUT_MS))
            != pdTRUE) {
        return ESP_ERR_TIMEOUT;
    }

    /* Find and unlink */
    for (star_bus_config_t* curr = manager->buses; curr;
         prev = curr, curr = curr->next) {
        if (curr->name && strcmp(curr->name, name) == 0) {
            to_remove = curr;
            if (prev == NULL) {
                manager->buses = curr->next;
            } else {
                prev->next = curr->next;
            }
            to_remove->next = NULL;
            ret = ESP_OK;
            break;
        }
    }

    if (ret == ESP_ERR_NOT_FOUND) {
        goto remove_give_mutex;
    }

    /* Cleanup operations */
    esp_err_t destroy_result = star_bus_config_destroy(to_remove);
    if (destroy_result != ESP_OK) {
        ret = destroy_result;
    }

remove_give_mutex:
    xSemaphoreGive(manager->mutex);
    return ret;
}
\end{lstlisting}

\subsubsection{Factory Function Pattern}

\begin{lstlisting}[language=C]
star_bus_config_t* star_bus_config_create_i2c(const char* name,
                                               i2c_port_t  port,
                                               uint8_t     address,
                                               gpio_num_t  sda_pin,
                                               gpio_num_t  scl_pin,
                                               uint32_t    clk_speed)
{
    if (name == NULL || strlen(name) == 0) {
        ESP_LOGE(s_TAG, "Invalid name for I2C bus config");
        return NULL;
    }

    star_bus_config_t* config = calloc(1, sizeof(star_bus_config_t));
    if (config == NULL) {
        ESP_LOGE(s_TAG, "Failed to allocate memory for I2C config");
        return NULL;
    }

    /* Initialize common fields */
    config->name        = strdup(name);
    config->type        = k_star_bus_type_i2c;
    config->initialized = false;
    config->handle      = NULL;
    config->next        = NULL;

    /* Initialize I2C-specific fields */
    config->proto.i2c.port    = port;
    config->proto.i2c.address = address;
    config->proto.i2c.config.sda_io_num = sda_pin;
    config->proto.i2c.config.scl_io_num = scl_pin;
    config->proto.i2c.config.master.clk_speed = clk_speed;

    /* Set default operations */
    star_bus_i2c_init_default_ops(&config->proto.i2c.ops);

    return config;
}
\end{lstlisting}

\subsection{Callback Functions}

\begin{lstlisting}[language=C]
/* Callback type definition */
typedef esp_err_t (*star_bus_callback_t)(star_bus_config_t* bus,
                                          void*              user_ctx);

/* Function using callback */
esp_err_t star_bus_manager_with_bus(star_bus_manager_t* manager,
                                     const char*         name,
                                     star_bus_callback_t callback,
                                     void*               user_ctx)
{
    ESP_RETURN_ON_FALSE(manager && name && callback,
                        ESP_ERR_INVALID_ARG, s_TAG, "Invalid arguments");

    /* ... find bus and invoke callback while holding mutex ... */
    esp_err_t callback_result = callback(found_bus, user_ctx);

    xSemaphoreGive(manager->mutex);
    return callback_result;
}

/* Using the callback pattern */
esp_err_t my_bus_callback(star_bus_config_t* bus, void* ctx) {
    ESP_LOGI(TAG, "Bus type: %s", star_bus_type_to_string(bus->type));
    return ESP_OK;
}

star_bus_manager_with_bus(&manager, "sensor_i2c", my_bus_callback, NULL);
\end{lstlisting}

\subsection{Summary}

\begin{itemize}
    \item Action functions return \texttt{esp\_err\_t}; predicates return \texttt{bool}
    \item Use \texttt{static} with \texttt{internal\_} prefix for file-local functions
    \item Use \texttt{priv\_} prefix for exposed-but-not-public functions
    \item Validate all parameters at function start
    \item Align parameters vertically in multi-line declarations
    \item Document all public functions with complete Doxygen
    \item Use single exit point with cleanup for complex functions
    \item Use factory functions for object creation
\end{itemize}

